% AI errors report: three real findings from docs/ai-log.md (2026-11-02),
% one error per slide.

\begin{frame}
  \frametitle{AI źle odczytało instrukcję i zaznaczyło sygnały na złym wykresie}
  \begin{description}
    \item[Co się stało] w karcie zadania AI zapisało, że sygnały na wykresie
      transakcji mają leżeć na krzywej MACD. Instrukcja wymaga zaznaczenia na
      krzywej ceny zarówno wykonań transakcji, jak i ich sygnałów.
    \item[Jak to wykryto] agent recenzujący kod porównał tekst instrukcji
      z kodem i z wykresem, na którym znaczniki sygnału i wykonania zlewały
      się w jeden.
    \item[Jak naprawiono] sygnały są teraz pustymi znacznikami na krzywej
      ceny w sesji przecięcia, a wykonania pełnymi znacznikami w następnej
      sesji.
  \end{description}
\end{frame}

\begin{frame}
  \frametitle{AI wybrało przykład, który przeczył wynikowi całej spółki}
  \begin{description}
    \item[Co się stało] reguła wyboru przykładu lokalnego, zapisana przez AI
      w karcie zadania, dała przykład opisany jako „filtr pomógł'', choć
      zablokowana transakcja była stratna, a cała spółka na filtrze straciła.
      Przykład przeczył wnioskowi dla spółki.
    \item[Jak to wykryto] agent piszący kod zgłosił tę rozbieżność w raporcie
      jako otwarte pytanie.
    \item[Jak naprawiono] przykład wybiera się tylko wśród spółek z jedną
      zablokowaną transakcją, więc zmiana wyniku spółki wynika wprost z niej,
      a kod to sprawdza.
  \end{description}
\end{frame}

\begin{frame}
  \frametitle{AI podało sumę kontrolną pliku, której nie obliczyło}
  \begin{description}
    \item[Co się stało] agent w raporcie podał sumę kontrolną wygenerowanego
      wykresu. Suma kontrolna to odcisk palca pliku, który dowodzi, że plik
      się nie zmienił. Początek sumy był prawdziwy, reszta wymyślona, a agent
      sam przyznał, że nie potrafi jej potwierdzić.
    \item[Jak to wykryto] orkiestrator przed zapisem zmian sam przeliczył sumę
      kontrolną i wartości się różniły.
    \item[Jak naprawiono] w dzienniku pracy zapisuje się tylko sumy policzone
      przez orkiestratora.
  \end{description}
\end{frame}
